
\begin{document}

\title{\textbf{Turing Machine Converter}}
\author{William Gregory, Matthew Wanta, and Dr. Ryan Dougherty}
\date{}
\maketitle
\thispagestyle{fancy}

The following is an auto-explainer for Turing Machines (TM) as represented by a series of figures depicting the tape as it is moving. \\

A computation of $T = (Q, \Sigma, \Gamma, \delta, q_0, h_a, h_r)$ on input string w is a sequence of configu-
rations, where:
\begin{itemize}
\item The first configuration is $q_0w$.
\item Each configuration yields the next via a single transition. To observe what character to read,
the TM looks at the cell to the right of where the state name is.
\item $T$ is in state $q$ observing symbol $y_1$.
\end{itemize}
If the transition moves L in the leftmost cell on the tape, the TM automatically stops and rejects.

\section{Turing Machine Definition}

A Turing Machine is defined as a 7-tuple 
\[
(Q, \Sigma, \Gamma, \delta, q_0, h_a, h_r)
\]
where:
\begin{itemize}
    \item $Q$ is a finite set of states,
    \item $\Sigma$ is a finite input alphabet,
    \item $\Gamma$ is a finite tape alphabet with $\Sigma \subset \Gamma$, and $\sqcup \in \Gamma \setminus \Sigma$ is the blank symbol,
    \item $q_0 \in Q$ is the start state,
    \item $h_a, h_r \in Q$ are the accept and reject states respectively (with $h_a \neq h_r$),
    \item $\delta$ is the transition function.
\end{itemize}
